% ---
% title: Defaults in a small loan portfolio
% tags: finance, binomial, expectation, variance
% ---

\begin{problem}
	\begin{context}
		A bank holds a small portfolio of loans to similar firms. A loan
		\emph{defaults} if the borrower does not pay it back. To keep the model
		simple, we assume that the borrowers default independently of each
		other: one firm's default does not change the chance that another firm
		defaults.
	\end{context}
	A bank has given $10$ loans. Each loan defaults within a year with
	probability $0.05$, independently of the others. Let $X$ be the number of
	defaults within the year.
	(a) Which distribution does $X$ have, with which parameters?
	(b) Find the probability that no loan defaults.
	(c) Find the probability that exactly one loan defaults.
	(d) Find the probability that at least two loans default.
	(e) Find $\E[X]$ and $\SD(X)$.
\end{problem}

\begin{slides}
	\small
	\exstatement
	\pause
	\begin{columns}[T]
		\begin{column}{.62\textwidth}
			\sol (a) $10$ independent loans, each defaults with the same
			probability $0.05$: \alert{binomial with $(n,p)=(10,\,0.05)$},
			$P(X=i)=\binom{10}{i}\,0.05^i\,0.95^{10-i}$.\\
			\pause
			(b) $P(X=0)=0.95^{10}\approx\alert{0.5987}$\\
			(c) $P(X=1)=10\cdot0.05\cdot0.95^{9}\approx\alert{0.3151}$\\
			\pause
			(d) Complement of $\{X=0\}\cup\{X=1\}$ (2 terms instead of 9):
			$P(X\ge2)=1-P(X=0)-P(X=1)\approx\alert{0.0861}$\\
			\pause
			(e) $\E[X]=np=\alert{0.5}$, \
			$\SD(X)=\sqrt{np(1-p)}=\sqrt{0.475}\approx\alert{0.689}$
		\end{column}
		\begin{column}{.36\textwidth}
			\centering
			\begin{tikzpicture}[x=.55cm,y=3.4cm]
			\draw[gray!40,very thin] (-0.6,0.2)--(5.6,0.2);
			\node[left,font=\tiny] at (-0.6,0.2) {$0.2$};
			\draw[gray!40,very thin] (-0.6,0.4)--(5.6,0.4);
			\node[left,font=\tiny] at (-0.6,0.4) {$0.4$};
			\draw[gray!40,very thin] (-0.6,0.6)--(5.6,0.6);
			\node[left,font=\tiny] at (-0.6,0.6) {$0.6$};
			\draw[->] (-0.6,0)--(5.8999999999999995,0) node[right,font=\scriptsize] {$i$};
			\draw[->] (-0.6,0)--(-0.6,0.75) node[above,font=\scriptsize] {$P(X=i)$};
			\filldraw[fill=gold!55,draw=gold] (0-.38,0) rectangle (0+.38,0.5987);
			\node[below,font=\tiny] at (0,0) {$0$};
			\node[above,font=\tiny] at (0,0.5987) {(b)};
			\filldraw[fill=gold!55,draw=gold] (1-.38,0) rectangle (1+.38,0.3151);
			\node[below,font=\tiny] at (1,0) {$1$};
			\node[above,font=\tiny] at (1,0.3151) {(c)};
			\filldraw[fill=teal!22,draw=teal] (2-.38,0) rectangle (2+.38,0.0746);
			\node[below,font=\tiny] at (2,0) {$2$};
			\filldraw[fill=teal!22,draw=teal] (3-.38,0) rectangle (3+.38,0.0105);
			\node[below,font=\tiny] at (3,0) {$3$};
			\filldraw[fill=teal!22,draw=teal] (4-.38,0) rectangle (4+.38,0.001);
			\node[below,font=\tiny] at (4,0) {$4$};
			\node[below,font=\tiny] at (5,0) {$5$};
			\draw[navy,thick,decorate,decoration={brace,amplitude=3pt}] (1.62,.13)--(5.4,.13) node[midway,above=3pt,font=\tiny] {(d) $X\ge2$: $0.086$};
			\draw[navy,thick,dashed] (0.5,0)--(0.5,.7);
			\node[right,font=\tiny,text=navy] at (0.5,.7) {$\E[X]=0.5$};
			\end{tikzpicture}
		\end{column}
	\end{columns}
\end{slides}

\begin{solution}
	\idea Each loan is one trial that either ends in default or not. The
	trials are independent and have the same default probability, so the
	number of defaults is binomial. The probabilities come from the binomial
	formula, ``at least two'' from the complement, and $\E[X]$, $\Var(X)$ from
	the formulas $np$ and $np(1-p)$.

	\step{(a)} Think of each of the $10$ loans as a trial, and call a default
	a ``success''. The trials are independent (the borrowers default
	independently), and every trial has the same success probability
	$p=0.05$. The number of successes in $n$ such trials is binomial with
	parameters $(n,p)$. So $X$ is binomial with parameters $(10,\,0.05)$:
	\[ P(X=i)=\binom{10}{i}\,0.05^i\,0.95^{10-i},\qquad i=0,1,\dots,10. \]

	\step{(b)} No default means all $10$ loans are repaid. Put $i=0$ in the
	formula; $\binom{10}{0}=1$ and $0.05^0=1$:
	\[ P(X=0)=0.95^{10}\approx0.5987. \]
	This also follows directly from independence: each loan is repaid with
	probability $0.95$, and we multiply these ten probabilities.

	\step{(c)} Put $i=1$. The factor $\binom{10}{1}=10$ counts which of the
	$10$ loans is the one that defaults:
	\[ P(X=1)=10\cdot0.05\cdot0.95^{9}\approx10\cdot0.05\cdot0.6302
		\approx0.3151. \]

	\step{(d)} ``At least two'' means $X=2,3,\dots,10$: nine terms. The
	opposite event is ``at most one'', $X=0$ or $X=1$, only two terms, and we
	already have both. The events $\{X=0\}$ and $\{X=1\}$ are mutually
	exclusive, so their probabilities add. Hence
	\[ P(X\ge2)=1-P(X\le1)=1-P(X=0)-P(X=1)\approx1-0.598737-0.315125
		\approx0.0861. \]
	(Subtracting the rounded values $0.5987$ and $0.3151$ would give $0.0862$;
	keep more digits in intermediate results.)

	\step{(e)} For a binomial random variable with parameters $(n,p)$,
	$\E[X]=np$ and $\Var(X)=np(1-p)$. Here
	\[ \E[X]=10\cdot0.05=0.5,\qquad \Var(X)=10\cdot0.05\cdot0.95=0.475,\qquad
		\SD(X)=\sqrt{0.475}\approx0.689. \]
	On average the bank expects half a default per year; the number of
	defaults typically differs from this by about $0.7$.
	\begin{answer}
		(a) binomial with parameters $(10,\,0.05)$ \quad
		(b) $0.95^{10}\approx0.5987$ \quad
		(c) $10\cdot0.05\cdot0.95^9\approx0.3151$ \quad
		(d) $\approx0.0861$ \quad
		(e) $\E[X]=0.5$, $\Var(X)=0.475$, $\SD(X)\approx0.689$
	\end{answer}
	\begin{remark}
		In (d) the complement saves work: two terms, $P(X=0)$ and $P(X=1)$,
		instead of the nine terms $P(X=2),\dots,P(X=10)$. Use it whenever an
		``at least'' question has a short opposite.
		\textbf{Common mistake:} forgetting the factor $\binom{10}{1}=10$ in (c).
		The probability $0.05\cdot0.95^9$ is the chance that one
		\emph{particular} loan defaults and the other nine do not; any of the
		ten loans can be the one.
	\end{remark}
\end{solution}
